\documentclass[11pt]{article}
\usepackage[margin=1in]{geometry}
\usepackage{amsmath,amssymb,amsthm,mathtools,booktabs,hyperref}
\newtheorem{definition}{Definition}
\newtheorem{theorem}{Theorem}
\newtheorem{lemma}{Lemma}
\newtheorem{remark}{Remark}
\title{A Hybrid Sustainability-Systems Architecture:\\Material Accounting, Information, Institutions, and Governability}
\author{Author(s) redacted for review}
\date{}
\begin{document}\maketitle

\begin{abstract}
Sustainability claims concern materially embodied systems whose physical dynamics, ecological function, observations, institutional authority, implementation, and normative constraints are jointly limiting. We present an axiomatic hybrid architecture that distinguishes material feasibility, ideal full-information viability, epistemic-output-feedback viability, and safety under a specified institution. The physical state is a history; the institutional decision state is a compatible-state set together with an institutional mode. Material balance is boundary-relative and includes continuous flows and event jumps. The framework does not claim a universal ecological law, thermodynamic closure, a general output-feedback safety theorem, or a universal policy-delay threshold. Instead it provides typed interfaces, quantifier conventions, nested safety objects, and admission standards for domain modules. The central scientific hypothesis is that sustainability failure can arise when a physically viable system lacks a common safe action under incomplete information, limited authority, or unreliable implementation.
\end{abstract}

\section{Status and scope}
\paragraph{Status convention.} Definitions describe the framework. Theorems and lemmas are conditional on their stated hypotheses. Templates and conjectures are not established results. This paper supplies architecture; its companion mathematical paper develops conditional kernels and proof obligations.

Material accounting is relative to a declared boundary. The represented material vector contains all represented mass-bearing compartments inside that boundary; omitted transfers must appear as declared boundary flows or structural discrepancy. Energy/exergy is an external driver in this paper. Material closure does not imply energetic feasibility, unlimited recyclability, or thermodynamic sustainability.

\section{Physical, functional, and institutional layers}
Let
\[
x=(r,f)\in\mathcal X_{\rm cont}:=\mathbb R_+^{n_r}\times\mathcal F,\qquad h\in\mathcal H.
\]
The vector $r$ contains material compartments; $f$ contains non-material ecological/functional states; and $h$ is an institutional mode/rule state that may be discrete. Let $\omega$ be external drivers, $\vartheta$ parameters, and
\[
\phi_t=x_t\in\mathscr H:=C([-\tau_{\max},0],\mathcal X_{\rm cont})
\]
the physical history.

Between a locally finite sequence of institutional or physical event times $t_k$,
\begin{align}
\dot r(t)&=\mathsf S\nu(\phi_t,h(t),u(t),\omega(t),\vartheta)+b(\phi_t,h(t),u(t),\omega(t),t),\label{eq:flow-r}\\
\dot f(t)&=G(\phi_t,h(t),u(t),\omega(t),\vartheta).\label{eq:flow-f}
\end{align}
At an event,
\begin{align}
r(t_k^+)&=\Delta_r(r(t_k^-),f(t_k^-),h(t_k^-),u_k,\xi_k),\\
f(t_k^+)&=\Delta_f(\cdot),\\
h(t_k^+)&=\Delta_h(h(t_k^-),Y_{[t_k-\tau_o,t_k]},\rho_k).
\end{align}
The post-event continuous history is defined in a declared hybrid or c\`adl\`ag phase space when $r$ or $f$ jumps; it is not forced into $C([-\tau_{\max},0],\mathcal X_{\rm cont})$ across a discontinuity. If $h$ is observed, the decision state is $(B_t,h_t)$ with $B_t$ conditional on that known mode. If $h$ is latent, it is included once in a joint belief and is not appended as a second coordinate.
Every module must state its event schedule, dwell-time rule, reset convention, priority rule, and implementation hold/interpolation rule.

\section{Boundary-relative material admissibility}
Elementary material fluxes obey $\nu_\ell\ge0$; reverse transformations use separate columns. Every material withdrawal from donor $r_j$ contains an instantaneous factor $\psi_j(r_j(t))$ satisfying $\psi_j(0)=0$. Negative boundary flows are also donor-limited.

Let $\mathsf L$ encode declared conserved moieties. If $\mathsf L^\top\mathsf S=0$, the hybrid material balance is
\begin{equation}
\mathsf L^\top r(t)-\mathsf L^\top r(0)
=
\int_0^t\mathsf L^\top b(\phi_s,u(s),\omega(s),s)\,ds
+
\sum_{t_k\le t}\mathsf L^\top[r(t_k^+)-r(t_k^-)].
\label{eq:hybrid-balance}
\end{equation}
An internal event is moiety-conserving exactly when its summand vanishes. A nonzero summand must be classified as a boundary exchange, measurement/reclassification event, or model discrepancy.

\paragraph{Material feasibility template.} A domain module must verify flow and reset positivity:
\[
\phi_j(0)=0\Rightarrow[\mathsf S\nu(\phi,\ldots)+b(\phi,\ldots)]_j\ge0,
\qquad
\Delta_r(\mathbb R_+^{n_r},\ldots)\subseteq\mathbb R_+^{n_r}.
\]
Positivity is not boundedness, persistence, recovery, or viability above thresholds.

\section{Safety, service, allocation, and normative constraints}
Services are typed outputs,
\[
s=\mathcal F(\phi,h,u,\omega,\vartheta).
\]
A state-safe set is
\[
\mathcal C_X(\lambda)=\{(\phi,h):g_i(\phi,h;\lambda)\ge0,\ i=1,\ldots,q\},
\]
where $\lambda\in\Lambda$ indexes declared normative thresholds, rights, service floors, and burden limits. A service/flow requirement is generally a state-control relation,
\[
\mathcal A(\phi,h,\omega;\lambda)=\{u\in\mathcal U(\phi,h):s(\phi,h,u,\omega,\vartheta)\ge s^{\min}(\lambda)\}.
\]
Distributional and procedural requirements enter the institutional authority set, rather than being inferred from aggregate feasibility.

Define a noncompensatory margin vector
\[
m(\phi,h,u;\lambda)=\bigl(g_1,\ldots,g_q,\ s_1-s_1^{\min},\ldots,s_p-s_p^{\min}\bigr)^\top.
\]
Safety requires $m\in\mathbb R_+^{q+p}$. A scalar certificate may be used only in noncompensatory form, e.g. $q(\phi,h,u)=\min_i m_i(\phi,h,u)$; weighted summaries may report preferences but do not certify componentwise safety without a restricted-domain proof.

\section{Observation, belief, authority, and implementation}
Institutions observe
\[
Y(t)=\mathcal O(\phi_t,h(t))+\varepsilon(t).
\]
At a decision time, the compatible-state set is
\[
B_t\subseteq\mathscr H\times\Theta,
\]
containing all histories and parameters consistent with declared initial uncertainty, actions, observations, errors, disturbance assumptions, and institutional records. The decision state is
\[
Z_t=(B_t,h_t).
\]

A prescription is not automatically an implemented physical control:
\[
a_t\in\Gamma(B_t,h_t),\qquad u_t\in\mathcal E(B_t,h_t,a_t),
\]
where $\Gamma$ captures authority, legal, budgetary, allocation, and procedural constraints, while $\mathcal E$ captures implementation, enforcement, and compliance uncertainty.

\paragraph{Quantifier convention.} Robust institutional safety uses the lower-game order
\[
\exists\text{ non-anticipative prescription strategy }a
\quad\forall\text{ compatible implementations, disturbances, parameters, and observation branches}.
\]
The alternative order $\forall w\exists u_w$ is not an implementable robust result unless the relevant $w$ is observed before action selection.

\section{Nested safety objects}
\begin{definition}[Material feasibility]
A history is materially feasible if the represented material stocks remain nonnegative under the declared flow, jump, and boundary rules.
\end{definition}

\begin{definition}[Full-information robust viability]
$\operatorname{RViab}^{\rm full}_T$ is the set of physical histories for which a non-anticipative full-state strategy keeps $X(t)\in\mathcal C_X$ and $u(t)\in\mathcal A(X(t),\omega(t);\lambda)$ for all $t\le T$ and all declared adversarial uncertainties.
\end{definition}

\begin{definition}[Epistemic-institutional viability]
Let $\mathfrak S$ be the set of admissible information states $(B,h)$. The kernel $\operatorname{IViab}_T$ contains information states from which there exists an institutionally authorized prescription strategy such that every compatible implementation, disturbance, parameter realization, and observation branch preserves safety between events through $T$.
\end{definition}

\begin{definition}[Institution-specific safety]
For a fixed institutional rule $\Pi$, $\operatorname{Safe}^{\Pi}_T$ contains initial information states from which the rule preserves safety against the declared uncertainty/implementation class.
\end{definition}

These objects are nested only when their information, authority, implementation, and uncertainty classes are aligned. Their differences diagnose distinct failures: physical, epistemic/common-action, authority, implementation, temporal, or recovery failure.

\paragraph{Failure taxonomy.} A sustainability assessment must distinguish: material inconsistency (unexplained moiety residual); physical infeasibility (no inward action at a physical boundary); epistemic infeasibility (no common safe action over $B$); authority infeasibility (a safe action exists but lies outside $\Gamma$); implementation infeasibility (the correspondence $\mathcal E$ admits unsafe realizations); temporal infeasibility (exit precedes informative review); recovery failure (no return to a viable information state); model-credibility failure (uncertainty/discrepancy coverage is unsupported); and normative incompatibility (declared state, service, or allocation constraints have empty intersection).

\section{Static reporting and information limits}
\begin{lemma}[Compensatory reporting limit]
For unrestricted margins $b\in\mathbb R^n$, $n\ge2$, and $w\in\mathbb R_{++}^n$, $w^\top b>0$ does not imply $b\in\mathbb R_+^n$.
\end{lemma}
\begin{lemma}[Static diagnostic aliasing]
If a safe and unsafe point state have the same instantaneous observation, no memoryless deterministic classifier of that observation correctly classifies both.
\end{lemma}
\begin{remark}
The latter lemma does not establish dynamic unobservability. Observer/filter/set-membership claims require dynamic observability/detectability, error, and structural-discrepancy conditions.
\end{remark}

\section{Modularity, information value, and recovery templates}
\paragraph{Compositional template.} Modules may exchange typed material/interface flows under assume--guarantee contracts. A compositional safety claim must verify interface routing, shared-control compatibility, and nonblocking hybrid-event composition; it cannot follow from separate local certificates alone.

\paragraph{Information template.} If information structure $\mathcal I_1$ causally refines $\mathcal I_2$ and authority/dynamics are fixed, finer information should not reduce viability because it may be ignored. This comparative statement becomes a conditional theorem only after belief dynamics and strategy spaces are specified.

\paragraph{Recovery template.} Recovery is not merely return of a physical state. It requires entry of the information/institution state $(B,h)$ into a viable kernel, possibly while remaining in a declared emergency envelope. Thus physical recoverability need not imply institutional recoverability.

\section{Substitution, spatial scope, and diagnostic scope}
Substitution is not assumed from aggregate capital or output. A claimed substitute must be represented as a typed pathway in $\mathcal F$, $\mathsf S\nu$, $\Gamma$, and $\mathcal E$, including its material inputs, energy requirement, delay, waste/displacement pathway, and authority/implementation conditions. This permits partial substitution without assuming that every ecological function has a technological replacement.

The framework is boundary-relative and can be lumped only when the claimed services and constraints are insensitive to spatial distribution. Where habitat configuration, hydraulic gradients, trade displacement, catchment routing, local exposure, or independent-unit fragility matters, a spatial/network module and aggregation rule are required. A stable aggregate stock or output is not itself evidence of local safety.

Component diagnostics are retained but typed by purpose. A gross throughput time is not a depletion horizon. A local net-rate horizon, a model-based boundary-crossing time, and a probabilistic failure horizon are distinct objects and belong in calibrated domain modules. No diagnostic alone establishes irreversibility unless recovery failure is part of the state model or supported by evidence.

\section{Admission standard and frontiers}
A domain module must supply: boundary and units; material and jump ledger; functional/service maps; state and state-control constraints; observation/proxy/error model; authority and implementation correspondences; uncertainty and discrepancy semantics; event model; identifiability plan; numerical verification; and out-of-sample/prospective validation.

Open frontiers include sampled epistemic-institutional fixed-point kernels, observer-to-safety transfer bounds, safe learning/active sensing, compositional certificates, hybrid nonlinear persistence, spatial/network extensions, and thermodynamic feasibility. These remain explicit because they are plausible routes to unification, not because they are established here.

\section{Conclusion}
Sustainability governability is posed as viability of a hybrid information-institution state constrained by material feasibility, ecological function, safety, authority, and implementation. The framework does not erase the gap between physical possibility and institutional realization; it makes that gap a central object for mathematical and empirical investigation.

\end{document}
